\setcounter{equation}{0}
\chapter{Introduction}
\sloppy

\section{Background}
Modern web browsers are used for banking, shopping, communication, education,
file exchange, and access to cloud applications. This convenience also makes the
browser an important attack surface. A user may encounter a phishing URL in a
message, a QR code that redirects to a counterfeit website, a malicious download,
a deceptive login form, an injected script, or a notification that imitates a
security warning. These attacks often depend on a sequence of browser events rather
than on one isolated file or address.

Traditional protection is distributed across browser warnings, blocklists, antivirus
software, email filters, and individual security extensions. These controls are
valuable, but they may not cover all browser activity, may provide little explanation,
and may not distinguish between a detected threat and an action that was actually
blocked. A URL may be new and absent from a blocklist, while a page may appear
normal until its scripts, frames, redirects, or forms are examined.

BrowseBlock proposes an additional, explainable protection layer. It combines
event-driven browser monitoring with machine-learning classification, static rules,
reputation evidence, webpage inspection, and authorized language-model-assisted
analysis. The design is intentionally modular so that each detector can be evaluated
independently and the user can see both the evidence and the limitations of a result.

\section{Problem Statement}
Users need timely protection against browser-based attacks without manually
investigating every URL, message, webpage, download, or browser extension. Existing
solutions are commonly specialized: URL classifiers may not analyze page behavior,
email filters may not follow extracted links, and heavyweight dynamic analysis is
unsuitable for continuous execution inside a browser extension. A blocklist can also
miss newly created or modified campaigns.

The problem addressed by BrowseBlock is the design of a unified security platform
that observes supported browser events, selects the relevant detector, combines
available evidence, explains the resulting risk, and recommends or applies a
protective action where browser capabilities permit. The system must also represent
incomplete analysis explicitly. A timeout, unsupported format, unavailable browser
event, or failed parser must not be reported as proof that an item is safe.

\section{Objectives of the Project}
The objectives are organized around the complete protection workflow.

\subsection{User Authentication and Access Control}
Provide sign-in and sign-out through Firebase Authentication, verify Firebase ID
tokens at the FastAPI backend, and restrict reports, analytics, configuration, and
other protected operations to authorized users and roles.

\subsection{Event-Driven Browser Monitoring}
Monitor supported navigation, webpage content, scripts, frames, redirects,
notifications, downloads, email content, and observable extension activity without
running every detector on every page. The triggering event should determine the
smallest relevant analysis workflow.

\subsection{Threat Detection and Classification}
Detect ten threat categories: URL and fake-QR phishing, XSS, credential harvesting,
HTML/JavaScript redirects, clickjacking, cryptojacking, malicious downloads, email
phishing, browser notification scams, and browser-extension abuse. Each result shall
be classified as Safe, Suspicious, or High-risk with supporting indicators.

\subsection{Evidence-Based Explanations}
Show the user why an item was flagged, including suspicious URL features, form
destinations, redirect evidence, script indicators, file characteristics, email
signals, notification content, or permission mismatches. The explanation shall
identify uncertainty and distinguish evidence from a guaranteed conclusion.

\subsection{Warnings and Protective Actions}
Warn the user and, where supported by browser APIs, restrict navigation, cancel a
download, or provide another configured response. The system shall record whether
the requested action was completed, unavailable, cancelled, or not applicable.

\subsection{Dashboard Analysis and Reporting}
Provide a React dashboard for manual URL, QR-code, email, and supported-file scans.
The dashboard shall display results, evidence, threat details, analytics, and
authorized history or reports stored through the backend.

\subsection{Privacy-Preserving Analysis}
Collect only the content required for the selected detector, never request real
passwords, avoid unnecessary retention of raw content, and keep model credentials,
database connections, and authentication secrets away from client code.

\section{Scope and Limitations}
\subsection{Scope}
The scope includes a Manifest V3 Chromium extension, a React and Tailwind CSS
dashboard, a Python FastAPI backend, Firebase Authentication, MongoDB application
data, detection services, and optional AI, RAG, LLM, and reputation integrations.
The system covers lightweight real-time checks as well as deeper user-authorized
dashboard analysis. It includes a common result format, risk explanation, logging,
reporting, and evaluation procedures.

\subsection{Limitations}
Protection is bounded by browser permissions, Chromium API visibility, network
availability, model quality, detector coverage, external-service availability, and
the evidence exposed by a page or file. A warning after navigation is not equivalent
to prevention before page execution. Static analysis cannot establish every runtime
behavior, and a low-risk result cannot guarantee that an item is harmless. The
prototype will therefore report its analysis boundary and will complement rather
than replace antivirus, browser security controls, and enterprise monitoring.

\section{Significance of the Study}
BrowseBlock connects detection, explanation, and user action in one workflow.
General users receive understandable warnings, while security analysts and
organizations can inspect evidence and reports. The project is significant because
it treats privacy, uncertainty, and enforcement status as first-class security
requirements rather than presenting a model score as an unconditional verdict.

\section{Project Overview}
The extension observes supported browser events and performs lightweight checks or
forwards validated input to FastAPI. The backend authenticates protected requests,
dispatches detector modules, combines model and rule evidence, stores authorized
scan records in MongoDB, and returns a consistent result. The dashboard uses the
same API to submit deliberate scans and review detailed reports.

\begin{center}
\textit{Browser event $\rightarrow$ Extension $\rightarrow$ FastAPI
$\rightarrow$ Detection service $\rightarrow$ Risk result $\rightarrow$
Warning, report, or supported protective action}
\end{center}

\section{Expected Outcomes}
The Phase I outcome is a detailed requirements and design foundation for a modular
BrowseBlock prototype. It defines threat modules, interfaces, data entities,
authentication boundaries, privacy constraints, and evaluation metrics. Phase II
will implement selected modules, train and validate models on representative data,
test browser compatibility, measure precision, recall, false positives, latency,
resource usage, and verify whether advertised warning, blocking, cancellation,
quarantine, and logging actions were actually completed.

\section{Sustainable Development Goals}
BrowseBlock supports the Sustainable Development Goals (SDGs) by applying
responsible computing to a practical cybersecurity problem. The project is
primarily connected with SDG 4, SDG 8, SDG 9, SDG 16, and SDG 17. These
connections are not claims that the software independently achieves an SDG;
rather, they describe how the project contributes to the knowledge, safety,
innovation, and responsible digital practices promoted by the goals.

\subsection{SDG 4: Quality Education}
BrowseBlock supports quality education by helping students and other learners
use online learning platforms, digital libraries, email, and educational
resources more safely. The extension provides understandable warnings and
evidence instead of presenting only a technical model score. This can improve
security awareness and help users recognize phishing, malicious redirects,
unsafe downloads, deceptive notifications, and suspicious browser extensions.
The project also provides an educational example of how machine learning,
web technologies, privacy requirements, and secure software design can be
combined in a real-world application.

\subsection{SDG 8: Decent Work and Economic Growth}
Unsafe links, credential theft, malicious files, and browser-based scams can
cause financial loss and interrupt the work of individuals and small
organizations. By identifying suspicious activity early and recommending
appropriate protective actions, BrowseBlock can reduce avoidable exposure
during online work. Its lightweight extension-based workflow is intended to
support secure browsing without requiring expensive enterprise security
infrastructure, making the approach relevant to users and organizations with
limited resources.

\subsection{SDG 9: Industry, Innovation and Infrastructure}
The project contributes to SDG 9 through an innovative and modular security
architecture. A Manifest V3 extension, FastAPI services, machine-learning
detectors, reputation checks, a React dashboard, Firebase Authentication, and
MongoDB are connected through defined interfaces. This architecture supports
incremental improvement: new detectors, evidence sources, or evaluation
methods can be added without redesigning the complete platform. The project
also encourages reliable digital infrastructure by treating availability,
latency, resource limits, and incomplete analysis as explicit design concerns.

\subsection{SDG 16: Peace, Justice and Strong Institutions}
Digital safety and trust are important components of inclusive and accountable
institutions. BrowseBlock supports this goal by protecting users from common
online deception and by recording evidence, risk classifications, explanations,
and protective-action status in an auditable result model. Authentication,
authorization, privacy controls, and clear reporting boundaries help prevent
unauthorized access to scan history and reports. The system is designed to
communicate uncertainty honestly, so that a failed or incomplete analysis is
not incorrectly represented as proof that content is safe.

\subsection{SDG 17: Partnerships for the Goals}
Effective cybersecurity depends on cooperation among developers, researchers,
educators, users, institutions, and providers of threat intelligence or
reputation services. BrowseBlock is designed with integration boundaries for
external data sources, authorized AI-assisted analysis, model evaluation, and
future collaboration. The project documentation, common result format, and
explicit detector limitations can support responsible sharing of findings and
future development while protecting credentials and unnecessary user content.
